% =================================================================================
% Response to the Editor (meta-review).
%
% NUMBERING IS BY ORDER OF APPEARANCE: each \begin{editorcomment} steps the
% editor counter, so the Nth editorcomment in this file becomes ``Editor
% Comment N''.  Move the environments and the numbers move with them.  Each
% editorcomment should be followed immediately by its editorresponse.
%
% The editor's summary is a separate spec from the reviewers' reports -- it is
% what the person who decides the outcome chose to emphasise.  Answer it on its
% own terms even when every point is repeated below.
% =================================================================================
\section{Response to the Editor}

% Paste the decision letter's editorial paragraph verbatim.  The environment
% already sets it in italic -- do not add quotation marks or \emph.
\begin{editorcomment}
{{EDITOR-COMMENT-VERBATIM}}
\end{editorcomment}

\begin{editorresponse}
We thank you for handling our submission and for the constructive feedback from the review panel. We have revised the manuscript following every comment from {{REVIEWER-LIST}}, and we summarize the most substantial changes below.

% One enumerate item per high-level change.  Same shape as the Summary of
% Revisions in main.tex: bold headline, concrete substance, then the comment
% ids the change answers.  Keep this list consistent with that one -- a
% mismatch between the editor summary and the point-by-point body is the first
% thing a re-reviewer notices.
\begin{enumerate}[leftmargin=*, topsep=2pt, itemsep=2pt]
    \item \textbf{{{CHANGE-HEADLINE}}} {{CHANGE-DESCRIPTION-WITH-CONCRETE-NUMBERS}} Addresses {{COMMENT-IDS}}.

    \item \textbf{{{CHANGE-HEADLINE}}} {{CHANGE-DESCRIPTION-WITH-CONCRETE-NUMBERS}} Addresses {{COMMENT-IDS}}.
\end{enumerate}

We detail each comment and our response in the following sections, preserving the original reviewer numbering from the decision letter.
\end{editorresponse}
